\scanpage{10}
l'unique autre point de $I$, à part $x_{-1},x_0,x_1,x_2$, qui sont reliés à $u$). Donc \deleted{\unsure{l'ens|l'un}} \inserted{chacun} des points \deleted{fixes} de la calotte standard \deleted{\unsure{est}} de centre $u$ est fixe. La $=$ chose est vrai pour les points de la calotte de centre $x_0$, donc aussi [appliquant \unsure{ceci|ces} aux triangles $(u,x_i,x_{i+1})$) à ceux \deleted{\unsure{la}} calottes \unsure{de|du} centre $x_i$, $i\in\mathbb Z/5\mathbb Z$. Mais Cela donne \uline{tous} les points \inserted{de $I$} sauf $u'$, or celui-ci doit être fixe aussi comme le complémentaire du \unsure{V|v|l'} ens précédent !

\deleted{Ainsi, on} On peut aussi regarder l'ens des faces et l'ens des arrêtes de l'icosaèdre $I$ \inserted{et leur relation d'incidence} (NB une arrête est déterminée par l'ens. des deux faces qui la contiennent), ils forment encore un graphe (dont les ``sommets'' sont les ``faces de l'icosaèdre ; \deleted{\unsure{en fait|on}} \inserted{ce} serait le graphe des sommets-arrêtes du dodécaèdre, \uline{dual} à l'icosaèdre). Je dis que le groupe des automorphismes de ce graphe est encore $\unsure{\Gamma|\Gamma'}$ [NB il est \deleted{immédiat} \inserted{évident} que $\Gamma$ y opère de façon fidèle]. On est encore ramené \unsure{à} prouver que tout autom. de ce graphe qui \deleted{\unsure{laisse}} invarie une face donnée et ses trois arrêtes, disons \deleted{\unsure{u,x_0,x_1}} encore $(u,x_0,x_1)$, est l'identité. Or il invarie alors chacun des \inserted{trois} triangles adjacents (comme les uniques triangles adjacents distincts sur les trois arrêtes envisagées), \deleted{et \unsure{on sait|invarie} alors leurs arrêtes (?)} puis les \deleted{triangles} faces $(u,x_2,x_3)$ et $(u,x_3,x_4)$ comme les seul couple de faces adjacent mutuellement, \uline{et} adjacents resp. à $(u,x_1,x_2)$ et $(u,x_0,x_{-1})$ (dont on sait qu'ils sont invariants). Cela donne l'invariance des \deleted{\unsure{les faces des deux}} arrêtes $(u,x_2)$ et $(u,x_4)$ commes intersection, donc aussi des arrêtes restantes $(x_1,x_2)$ et $(x_0,x_{-1})$ des deux \deleted{triangles} faces $(u,x_1,x_2),(u,x_0,x_{-1})$. Donc leurs arrêtes sont fixes, \deleted{\unsure{\illegible{2}|donc}} et par symétrie on voit que pour chacune des \deleted{triangles} \inserted{faces} adjacentes à $(u,x_0,x_1)$, \deleted{les \unsure{trois}} ces faces \uline{et} leurs arrêtes sont fixes. À partir de là la conclusion est triviale
% archive sheet number: 9
